\subsection{Sprint 1}

    Para esta sprint estavam previstas 6 User Stories, separadas em pequenas
tasks de Back e Front. Faziam parte das atividades previstas, as relacionadas aos
domínios de Usuário (login), Veículos, Motoristas, Passageiros dependentes e
outros, além de telas auxiliares para integração e utilitárias, como a Splash Screen,
por exemplo. Da minha parte estava previsto monitorar o Backend, ajudar os
colegas a clonarem e desenvolverem no repositório e revisar os Pull Requests
referentes ao Backend para garantir que a qualidade do código correspondesse ao
esperado. Além disso, tinha em mente entender melhor o deploy se houvesse algum
tempo ocioso. Não fazia parte do meu escopo previsto fazer tasks sozinho, então
realmente só as faria se precisasse.

    Tanto a equipe de Back quanto Front entregaram todas as User Stories
previstas, um número bem alto de componentes e endpoints, que até deixaram
nossos stakeholders (também estudantes de Engenharia de Software da PUCRS)
surpresos, tendo em vista que realmente puxamos muitas coisas nesta primeira
sprint. Da minha parte o Backend evoluiu muito bem pela metade da sprint, o
começo estava meio lento e, conforme requisitado por colegas Ages 4, eu e meu
colega Ages 3 do Backend pegamos tasks para fazer, o que junto à crescente da
equipe e a facilidade de validar código com a técnica de Test Driven Development
adotada, possibilitou que entregássemos o Backend inteiro pela metade da sprint,
sendo assim tirei um tempo para estudar Infra e diferentes serviços da AWS, além
de estudar a viabilidade e dificuldade de subir um Terraform.

    Chegando ao final da sprint, descobrimos um problema de histórico corrompido das migrations que
assustou o time; no entanto, como já havia resolvido isso no trabalho, fiz questão de
tranquilizar a equipe, assumir a responsabilidade de consertar e simplificar ao
máximo os comandos para que a equipe inteira conseguisse recriar o banco de
dados. Para acabar a sprint, novamente a pedido de colegas Ages 4 dei uma ajuda
no Frontend, que tinha um escopo um pouco maior, entregando a Splash Screen e
uma integração da tela de adicionar, remover ou editar dependentes.

    Confesso que coloquei a mão no código mais do que gostaria e conversando
com colegas Ages 3 e Ages 4, o sentimento também foi parecido. Sendo assim,
fizemos questão de reforçar isso na retrospectiva, até porque a situação foi pior
ainda para os Ages 3 do Front e os Ages 4 que se comprometeram a deixar a
entrega bem polida e portanto tiveram que se estender no horário um dia antes da
entrega.

    Creio que, em vez de lições concretas, tiro hipóteses desta sprint. A primeira é
que talvez a sprint tenha sido grande demais para Ages 2 e principalmente Ages 1, que
ainda estão meio perdidos no ecossistema da Ages e portanto não produzem com a
mesma velocidade que vemos de nós mesmos. A segunda é que talvez tenha faltado
alguém criar uma aproximação com alguns colegas mais novatos para que eles se
sentissem mais à vontade de participar e perguntar, facilitando o processo de
trabalho em equipe. Por fim, a última hipótese levantada é que, apesar de um
ambiente que permite amizades, a Ages deveria ser primeiramente de cunho
profissional. Sendo assim, apesar de acolhimento e amizade, é necessário cobrar
entregas mais rigidamente, coisa que vejo muitas vezes não acontecer, pelo fato da
equipe criar uma amizade que quebra um pouco da responsabilidade profissional
envolvida.

    Creio que está chegando a hora de termos um deploy concretizado, de
momento temos apenas um diagrama que ainda estamos validando entre nós
mesmos e com o arquiteto da Ages, no entanto, no mais tardar ao fim desta sprint,
acredito que devemos ter subido pelo menos o Backend, que menos trouxe
complicações, na AWS e até tentar esboçar um terraform para especificar este
deploy. Além disso planejo me aproximar cada vez mais de colegas de Ages,
principalmente Ages 1, no entanto para isso, creio que tenho que simplificar um
pouco mais o linguajar com que me comunico com eles e explicar com mais
paciência as coisas, para não assustá-los com uma complexidade que pode sim ser
abstraída para fins de compreensão e confiança deles.
